\documentclass[11pt]{article}

\usepackage[T1]{fontenc} 
\usepackage{geometry}
\usepackage{amsmath, amssymb, amsthm}
\usepackage{mhchem}
\usepackage{hyperref}

\usepackage{tikz}
\usetikzlibrary{calc}

\geometry{a4paper, margin=1in}

\theoremstyle{remark}
\newtheorem*{answer}{Answer}

\title{HU3101: History and Philosophy of Science}
\author{Satvik Saha}
\date{}

\begin{document}
    \noindent\textbf{IISER Kolkata} \hfill \textbf{Class Test VI}
    \vspace{3pt}
    \hrule
    \vspace{3pt}
    \begin{center}
    \LARGE{\textbf{HU3101: History and Philosophy of Science}}
    \end{center}
    \vspace{3pt}
    \hrule
    \vspace{3pt}
    Satvik Saha, \texttt{19MS154} \hfill \today
    \vspace{20pt}


    \paragraph{Question 1.} Describe the steps through which nuclear fission was
    discovered.

    \begin{answer} The discovery of nuclear fission is a key milestone in science,
    sandwiched between the discovery of the phenomenon of radioactivity, and the
    creation of the nuclear bomb.

        \begin{enumerate}
            \item \textit{Discovery of radioactivity:} X-rays were discovered by
            Wilhelm R\"ontgen in 1895. Attempting to produce these from uranium
            salts, Henri Becquerel ended up discovering an entirely new phenomenon
            merely a year later -- these salts were capable of producing an image on
            an X-ray plate without any external energy source. Marie and Pierre Curie
            further investigated this phenomenon, finding the same properties in
            thorium, polonium, and radium (the latter two being elements they
            discovered themselves). This was dubbed \emph{radioactivity}. 

            Ernest Rutherford classified these new rays into two types -- $\alpha$
            and $\beta$. Later, Paul Villard identified a third, namely $\gamma$. By
            1909, $\alpha$ particles had been identified as helium nuclei, and
            $\beta$ particles as electrons/cathode rays. Meanwhile, Rutherford and
            Curie both observed that these radioactive elements seemed to undergo
            decay, releasing a radioactive gas -- this was named radon. Rutherford
            and Frederick Soddy investigated these rates of decay of the radioactive
            products, and introduced the concept of a `half-life'. Later, Soddy and
            Margaret Todd introduced the idea of chemical isotopes -- atoms of the
            same element (chemically and spectroscopically identical) by having
            different half-lives.


            \item \textit{Investigations into the atomic nucleus:} In 1909,
            Rutherford performed his famous gold-foil experiment, firing $\alpha$
            particles at a thin sheet of gold. He deduced the existence of an atomic
            nucleus: a minuscule region within an atom which carried practically all
            of its mass, with the remainder being mostly empty space. The Rutherford
            model thus describes a nucleus containing positively charged protons,
            circled by negatively charged electrons. In 1913, Niels Bohr introduced
            the Bohr atomic model, a groundbreaking achievement.

            One problem with these models was that it was difficult to account for
            the stability of an atom, given the intense electrostatic repulsion
            between the positively charged protons in the nucleus. Rutherford and
            Bohr began hypothesising about another kind of nuclear particle, which
            would compensate for these repulsive forces by exerting an attractive
            nuclear force of its own. This was supported by a mismatch (by a factor
            of 2) between the observed atomic masses and the mass contributed by
            protons alone. In 1932, James Chadwick discovered and identified the
            neutron, confirming these suspicions.


            \item \textit{Nuclear transformations:} Following the discovery of
            radioactivity, a plethora of experiments were performed in which atoms
            were bombarded by these new $\alpha$ particles. One notable instance is
            that of the transmutation of nitrogen into oxygen, performed by
            Rutherford but actually identified by Patrick Blackett. \[
                \ce{^14_7N + ^4_2He \to ^17_8O + p^+}.
            \] James Chadwick's experiments used the following nuclear reaction. \[
                \ce{^9_4Be + ^4_2He \to ^12_6C + n^0}.
            \] It was noted that the electrically neutral neutron would likely serve
            as a much better bombarding particle. This idea was used by Enrico Fermi
            and his team in Rome, with help from Lise Meitner and Otto Hahn in
            Berlin. By performing neutron bombardment of elements such as aluminium
            and above, they were able to produce `induced radioactivity', i.e.\ the
            element would be converted into a radioactive isotope. Fermi identified
            three results of this neutron irradiation: release of $\alpha$ particles,
            release of a proton, or release of $\gamma$ rays. Finally, the newly
            created radioactive isotope would undergo $\beta$ decay, resulting in a
            transmutation of the element. In this way, Fermi generated what was
            thought to be new, heavier elements by bombarding uranium with neutrons,
            winning him the 1938 Nobel Prize in Physics.


            \item \textit{Nuclear splitting:} Fermi's discovery was met with
            skepticism by Ida Noddack, who pointed out that the neutron bombardment
            of uranium, a fairly heavy element, could have caused it to split into
            much lighter elements instead of neighbouring ones. This idea was not
            followed up upon, since there was currently no theoretical understanding
            of how a tiny neutron could make this happen. 

            By 1938, Meitner was forced to flee from Nazi Germany, and ended up in
            Sweden. She continued to work with Hahn and Fritz Strassman who were at
            Berlin. When redoing Fermi's experiments with uranium, they discovered
            that the products contained barium, a significantly lighter element than
            uranium. Thus, it was indeed the case that the bombarding neutrons
            cracked the uranium nucleus into sizeable chunks.

            Meitner discussed this puzzling result with her nephew Otto Frisch, and
            during the Christmas holiday, they went for a walk in the snow and made
            the following insights: using the liquid drop model of the nucleus
            advocated by Gamow and Bohr, they found a mechanism to explain how the
            incoming neutron could sufficiently deform the uranium nucleus, resulting
            in it splitting into two fragments. Now, these newly formed nuclei must
            repel each other and move away at very high speed; this would require
            around 200 MeV of energy. Where could this energy come from?

            Meitner calculated on the spot that the masses of the daughter nuclei
            empirically fall short of the mass of the original uranium nucleus by
            about one-fifth that of a proton. Enter Einstein's mass-energy
            equivalence principle, \[
                E = mc^2.
            \] Plugging in $m$ as the mass discrepancy gives precisely the required
            200 MeV of energy. The same was experimentally confirmed by Frisch in
            January 1939. After speaking with the biologist William Arnold, he named
            this process \emph{fission}, analogous to cell division. These results
            were published in \emph{Nature} soon afterwards.
        \end{enumerate}

        It was noted that this nuclear reaction resulted in the release of neutrons
        as well, which themselves can bombard other nuclei; this would lead to a
        \emph{chain reaction}. The relevant nuclear equation describing the fission
        of uranium-235 (one possibility) is as follows. \[
            \ce{n^0 + ^235_92U \to ^139_56Ba + ^94_36Kr + 3n^0}.
        \] There are two major paths of development after this: using such an
        uncontrolled chain reaction and its tremendous release of energy to create an
        atomic bomb, and harnessing a controlled chain reaction to create a power
        plant. The former was actively pursued during the Second World War, notably
        in the Manhattan Project. Two atomic bombs were dropped on the cities of
        Hiroshima and Nagasaki in 1945. The first nuclear reactor producing
        electricity began operating in 1951.

    \end{answer}

    \paragraph{Question 2.} Describe the four kinds of proof admitted in Indian
    philosophy, indicating the merits and demerits of each as valid scientific
    reasoning.

    \begin{answer}
        We focus on the Nyaya school of thought, which recognizes four
        \textit{pramanas}, i.e.\ means of knowledge.

        \begin{enumerate}
            \item \textit{Pratyaksha (perception):} This is perhaps the most
            important of the four; it is interesting to note that this is the
            \emph{only} method accepted by the Charavaka school of thought.

            Here, perception is divided into \emph{laukika} (ordinary) and
            \emph{alaukika} (extraordinary). Ordinary perception is simply the use of
            one's sense organs (sight, touch, taste, smell). Extraordinary perception
            includes \emph{pratibha} (intuition), \emph{samanyalaksanapratyaksa}
            (induction), and \emph{jnanalaksanapratyaksa} (retrograde analysis).

            In order for an observation to qualify as being valid, it must satisfy
            \emph{indriyarthasannikarsa} (direct experience), \emph{avyapadesya} (be
            non-verbal), \emph{avyabhicara} (be unchanging), and
            \emph{ayavasayatmaka} (be definite). Essentially, one must perceive an
            object directly, without relying on someone else. This observation must
            have some sort of stability, and must be free of doubt.

            \item \textit{Aumana (inference):} This is the use of previously known
            knowledge to derive new knowledge, by a combination of induction and
            deduction. A statement of inference will contain three parts -- minor,
            major, and middle. The \emph{paksha} or minor is the subject of
            inference, i.e.\ the thing we are trying to make an inference about. The
            \emph{sadhya} or major is the object of inference, i.e.\ the property of
            the subject which we are trying to establish. These are to be connected
            by the \emph{hetu} or reason/middle term, which is some general, true
            statement about the subject. Of course, the strength of such an argument
            is determined by the \emph{hetu}, which must be free of contradiction.
            The Nyaya darshan also identifies possible fallacies in the \emph{hetu},
            as follows.
            \begin{enumerate}
                \item \textit{Asiddha:} The \emph{hetu} is simply false/unproven. It
                may be the case that the subject is itself non-existent
                (\emph{ashrayasiddha}), which makes whatever follows meaningless.
                Another possibility is that the \emph{hetu} is not present in the
                subject at all (\emph{svarupasiddha}), or it is conditionally present
                (\emph{vyapyatvasiddha}).

                \item \textit{Savyabhichara:} The \emph{hetu} is not properly
                connected with the subject. For instance, the \emph{hetu} may be too
                general (\emph{sadharana}), too specific (\emph{asadharana}), or
                non-exclusive (\emph{anupasamhari}).

                \item \textit{Satpratipaksa:} The \emph{hetu} is contradicted by a
                different \emph{hetu}.

                \item \textit{Badhita:} The \emph{hetu} is contradicted by a
                superseding proof, such as perception.

                \item \textit{Viruddha:} The \emph{hetu} proves the opposite statement.
            \end{enumerate}

            \item \textit{Upamana (analogy):} This is the use of analogy/comparison
            with existing knowledge. More precisely, it is the knowledge of the
            relation between a name, and the object denoted by that name. This is
            often useful when talking about descriptions: comparing some object with
            a known one makes it easier to identify.

            \item \textit{Sabda (testimony):} This is the use of the word/testimony
            of others, presumably experts. It can be argued that one person cannot
            derive everything they need (using the previous means of knowledge)
            practically; this would require far too much time and energy. Thus, one
            has to rely on previously acquired knowledge, communicated by the spoken
            word or writing. \emph{Sabda} has to come from a trustworthy source;
            divine revelation, or other people. The former is said to be infallible,
            the latter is not.

        \end{enumerate}
        We note the absence of the pramanas \textit{Arthapatti} (derivation from
        circumstances) and \textit{Anupalabdi} (knowing a negative), which are
        sometimes accepted by other schools.

        In the context of modern scientific reasoning, we can reaffirm
        \emph{pratyaksha} as the most important source of knowledge. Science is
        almost entirely concerned with the observable/perceivable world, and when
        assessing a particular hypothesis or theory, raw observational data trumps
        everything else. Of course with advancements in technology, we have
        recognized that the human senses are not enough; thus, we use instruments to
        make up for these deficiencies, say a telescope or a microscope for better
        vision. This is also recognition that our senses can be unreliable, sometimes
        consistently so (optical illusions are a great way of demonstrating the
        unreliability of visual perception) which makes the use of these external
        tools even more important. The image below shows the Ebbinghaus illusion: the
        two orange circles are in actually identical. While this is difficult to
        directly perceive, we are perfectly happy to accept this fact once the
        circles are measured using a scale, such is our trust in these instruments.
        
        \begin{center}
        \begin{tikzpicture}[baseline=(X.base)]
            \node[circle,fill=orange,draw=orange,minimum size=1.5cm] (X) at (0,0) {};
            \foreach \i in {0,60,...,330}{
                \filldraw[blue!50!white]  (\i:2.8) circle (1.3);}
        \end{tikzpicture}
        \hspace{1cm}
        \begin{tikzpicture}[baseline=(X.base)]
            \node[circle,fill=orange,draw=orange,minimum size=1.5cm] (X) at (0,0) {};
            \foreach \i in {0,45,...,360}{
                \filldraw[blue!50!white]  (\i:1.2)  circle (.3);}
        \end{tikzpicture}
        \end{center}

        Similar arguments show that \emph{pratibha} or intuition is not a reliable
        method of reasoning; science today has brought many counter-intuitive results
        to light, an extreme example being quantum mechanics. Finally, it is not
        enough to focus on one's own perception, since modern scientific knowledge is
        based on repeatable, reproducible, testable results. Something perceived by a
        single individual amounts to anecdotal evidence; only when it has been tested
        by other independent individuals or bodies does it enter scientific knowledge
        (this seems to carry an element of \emph{sabda} or authority). These points
        are meant to highlight the untrustworthiness of perception when applied
        na\"ively under certain circumstances. 

        As we have seen, the use of \emph{anumana} or inference is heavily dependent
        on the \emph{hetu} or reason; thus, there has to be a great deal of
        confidence in this reason, which is typically a general principle applicable
        to the subject. Thus, the main problem becomes justifying the \emph{hetu},
        which ultimately comes down to the problem of induction (how can we guarantee
        that the reason applies to the subject, based on limited observations).
        Today, we use Popper's principle of falsifiability to establish this trust
        (rarely absolute) in the \emph{hetu}. Once this is done, there isn't really
        any doubt in the validity of the syllogisms that follow. Many objects in
        modern science, such as subatomic particles, are not directly observable (at
        least not in same sense as \emph{pratyaksha}), their existence can only be
        inferred from other observations. Similarly, one could also argue that any
        observation made by a scientific instrument is not perceived knowledge, but
        rather inferred knowledge (based on the assumption that the instrument does
        indeed work, and accurately reflects reality).

        \emph{Upamana} seems to be an essential part of discovery, specifically
        communicating that discovery to the rest of the world. When a novel phenomenon
        is first encountered, the job of the discoverer is to document it so that
        others can independently verify these findings. In this process, one has no
        choice but to use analogy; after all, the unknown has to be described by
        something already known for the description to make sense -- this is
        fundamental to human communication, which is based upon common experience. In
        modern science, this need only apply to the qualitative aspects of a
        description; the quantitative parts are specified using measurements. In the
        same way, \emph{upamana} is essential to scientific communication in general,
        as well as the learning process. Outside of this, \emph{upamana} does not
        play any role in science; it cannot be used to generate new knowledge.

        Like \emph{upamana}, \emph{sabda} or testimony is tied with scientific
        communication. A person today cannot independently verify (by
        \emph{pratyaksha} or \emph{anumana}) the vast amount of knowledge required to
        lead a normal life, there is simply too much. Thus, we rely on experts in
        various fields to communicate what they know to us. An important distinction
        between \emph{sabda} and its modern counterpart is that science is rooted in
        observations of nature, not divine revelation. Thus, in order for a piece of
        knowledge communicated by \emph{sabda} to be considered scientific, it must
        be possible to trace it back to a conventional source (something arising from
        the scientific method). Thus, \emph{sabda} is never the source (nor proof)
        of scientific knowledge, merely a tool of dissemination; \emph{sabda} used
        as an `argument from authority' without adequate evidence backing it is out
        of the question. In the presence of contradictory \emph{sabda}, there is no
        way for a casual listener (who presumably has no other way of gaining that
        knowledge) to judge with absolute certainty which is correct. Here, an
        element of trust creeps in, and it becomes necessary to talk about
        `scientific consensus'. Consensus among experts in a field can be taken as a
        heuristic, used to determine a level of trust.


    \end{answer}

\end{document}
